% This LaTeX document needs to be compiled with XeLaTeX.
\documentclass[10pt]{article}
\usepackage[utf8]{inputenc}
\usepackage{amsmath}
\usepackage{amsfonts}
\usepackage{amssymb}
\usepackage[version=4]{mhchem}
\usepackage{stmaryrd}
\usepackage{diagbox}
\usepackage[fallback]{xeCJK}
\usepackage{polyglossia}
\usepackage{fontspec}
\IfFontExistsTF{Noto Serif CJK TC}
{\setCJKmainfont{Noto Serif CJK TC}}
{\IfFontExistsTF{STSong}
  {\setCJKmainfont{STSong}}
  {\IfFontExistsTF{Droid Sans Fallback}
    {\setCJKmainfont{Droid Sans Fallback}}
    {\setCJKmainfont{SimSun}}
}}

\setmainlanguage{french}
\IfFontExistsTF{CMU Serif}
{\setmainfont{CMU Serif}}
{\IfFontExistsTF{DejaVu Sans}
  {\setmainfont{DejaVu Sans}}
  {\setmainfont{Georgia}}
}

\begin{document}
\section*{Analyse Numérique - correction TD n \({ }^{\circ} \mathbf{1}\) 05-02-2025 - SIGMA1}
\section*{Exercice 1:}
On a \(f(1 / 8)=\sin (\pi / 8)=0,3827\)\\
Les points d'interpolation sont \((0,0),(1 / 2,1)\) et \((1,0)\)\\
Le polynôme d'interpolation de degré 2 de Lagrange s'écrit :

\[
P_{2}(x)=0 \times \frac{\left(x-\frac{1}{2}\right)(x-1)}{\frac{1}{2}}+1 \times \frac{x(x-1)}{\frac{1}{2}\left(\frac{1}{2}-1\right)}+0 \times \frac{x\left(x-\frac{1}{2}\right)}{\frac{1}{2}}=-4 x(x-1)
\]

On a alors \(P_{2}(x)=-4 x(x-1)\) dans \([0,1]\) et \(P_{2}(0,125)=0,4375\)\\
Comme \(P_{21}\) est un polynôme de degré 2, on a alors besoin de 3 points.\\
Comme, on cherche à calculer \(P_{21}(1 / 8)\) et que \(1 / 8 \in[0,1 / 2]\), On va prendre l'intervalle [0, 1/2] pour construire notre polynôme d'interpolation de Lagrange.

On a les points d'interpolation \((0,0)\) et \((1 / 2,1)\) et on choisit le point du milieu de l'intervalle (1/4, \(\sqrt{ } 2\) /2). Le polynôme d'interpolation de Lagrange s'écrit alors :\\
\(P_{21}(x)=0 \times \frac{\left(x-\frac{1}{4}\right)\left(x-\frac{1}{2}\right)}{\frac{1}{8}}+\frac{\sqrt{2}}{2} \times \frac{x\left(x-\frac{1}{2}\right)}{\frac{1}{4}\left(\frac{1}{4}-\frac{1}{2}\right)}+1 \times \frac{x\left(x-\frac{1}{4}\right)}{\frac{1}{2}\left(\frac{1}{2}-\frac{1}{4}\right)}=-\frac{16 \sqrt{2}}{2} x\left(x-\frac{1}{2}\right)+8 x\left(x-\frac{1}{4}\right)\)\\
\(P_{21}(x)=-\frac{16 \sqrt{2}}{2} x\left(x-\frac{1}{2}\right)+8 x\left(x-\frac{1}{4}\right)\) dans \([0,1 / 2]\) et \(P_{21}(0,125)=0,4053\)

\section*{Exercice 2:}
Interpolation par des Polynômes de Lagrange de la fonction sinus.\\
On considère l'interpolation par Polynốmes de Lagrange de la fonction simus \((x)\), xe exprimé en radians, aux points \(\left(\frac{\pi}{2}, \frac{2 \pi}{3}, \frac{3 \pi}{4}, \frac{5 \pi}{6}\right)\).\\
La méthode des différences-divisées est utilisée.

\begin{enumerate}
  \item Trower le polynôme \(P(x)\) - Les grandeurs calculées sont représentees par heur valeure décimale.
  \item Estimer l'errew \(|P(x)-\sin (x)|\) lorsque
\end{enumerate}

\[
x \in\left[\frac{\pi}{2}, \frac{5 \pi}{6}\right]-
\]

\begin{enumerate}
  \setcounter{enumi}{2}
  \item Estimer l'eveneur \(|P(x)-\sin (x)|\) lorsque \(x=\pi\). Que de'duisez-vous sur he comporte. ment de l'erreur hors de l'intervalle d'interpolation?
  \item Le polynôme d'interpolation passant par les tepoints sera au maximum de degré 3.\\
Calculs interme'diaires
\end{enumerate}

\begin{itemize}
  \item \(f\left[x_{0}, x_{1}\right]=\frac{-0,1339745962 \times 6}{\pi}=\frac{-0,1339745962}{0,5235991}\)\\
\(f\left[x_{0}, x_{1}\right]=-0,2558724723\)
  \item \(f\left[x_{1}, x_{2}\right]=\frac{0.7071067-0.8660254038}{2.35619449-2.0943951}\)
\end{itemize}

\[
=\frac{-0,1589187038}{0,26179939}=-0,6070247291
\]

心

\[
\begin{aligned}
\cdot f\left[x_{2}, x_{3}\right] & =\frac{0,5-0,7071067}{2,617993876-2,35619449} \\
& =\frac{-0,2071067}{0,261799386}=-0,7910893267
\end{aligned}
\]

\begin{itemize}
  \item \(f\left[x_{0}, x_{1}, x_{2}\right]=\frac{-0,6070247291+0,2558724723}{2.35619449-1.570796}\)\\
\(=\frac{-0,3511522568}{0,78539849}=-0,4471007537\)
  \item \(f\left[x_{1}, x_{2}, x_{3}\right]=\frac{-0,7910893267+0,6070247291}{2,617993876-2,0943951}\)
\end{itemize}

\[
=\frac{-0,1840645976}{0,523598776}=-0,3515374864
\]

\begin{itemize}
  \item \(f\left[x_{0}, x_{1}, x_{2}, x_{3}\right]=\frac{-0.3515374864+0.4471007537}{2.617993876-1.570796}\)
\end{itemize}

\[
=+0,0912556896
\]

\begin{center}
\begin{tabular}{|l|l|l|l|l|}
\hline
\backslashbox{x}{\({ }^{\mathcal{C}}\) \&} & 0 \(P_{0}\) & \begin{tabular}{l}
1 \\
\(P_{1}\) \\
\end{tabular} & \begin{tabular}{l}
2 \\
\(P_{2}\) \\
\end{tabular} & \begin{tabular}{l}
3 \\
\(P_{3}\) \\
\end{tabular} \\
\hline
\begin{tabular}{l}
\(\frac{\pi}{2}=\) \\
1,5708 \\
\end{tabular} & 1 &  &  &  \\
\hline
\begin{tabular}{l}
\(\frac{2 \pi}{3}=\) \\
2.0944 \\
\end{tabular} & \begin{tabular}{l}
\(\frac{\sqrt{3}}{2}=\) \\
0,866 \\
025 \\
\end{tabular} & \begin{tabular}{l}
\(f\left[x_{0}, x_{1}\right]=\) \\
-0,25587247 \\
\end{tabular} & \multicolumn{2}{|c|}{} \\
\hline
\begin{tabular}{l}
\(3 \pi=\) \\
4 \\
2.3562 \\
\end{tabular} & \begin{tabular}{l}
\( \frac{\sqrt{2}}{2}= \) \\
0,7071 \\
067 \\
\end{tabular} & \begin{tabular}{l}
\( f\left[x_{1}, x_{2}\right]= \) \\
-0,6070247291 \\
\end{tabular} & \begin{tabular}{l}
\( f\left[x_{0}, x_{1}, x_{2}\right]= \) \\
-0,4471007537 \\
\end{tabular} &  \\
\hline
\begin{tabular}{l}
\( \frac{5 \pi}{6}= \) \\
261799 \\
\end{tabular} & \begin{tabular}{l}
\( \frac{1}{2} \) \\
0,5 \\
\end{tabular} & \begin{tabular}{l}
\( f\left[x_{2}, x_{3}\right]= \) \\
-0,7910893267 \\
\end{tabular} & \begin{tabular}{l}
\( f\left[x_{1}, x_{2}, x_{3}\right]= \) \\
-0,3515374864 \\
\end{tabular} & \begin{tabular}{l}
\( f\left[x_{0}, x_{1}, x_{2}, x_{3}\right]= \) \\
+0,0912556896 \\
\end{tabular} \\
\hline
\end{tabular}
\end{center}

Déù l'écriture du polynôme d'interpolation:\\
\(P_{0}(x)=1\)

\[
\begin{aligned}
P_{1}(x) & =P_{0}(x)+\left(x-x_{0}\right) f\left[x_{0}, x_{1}\right] \\
& =1+(x-1.570796)(-0,2558724723) \\
& =-0,2558724723 x+1+1,570796 \times 0,2558724 \ldots \\
& =-0,2558724723 x+1,40192
\end{aligned}
\]

\[
\begin{aligned}
& P_{2}(x)=P_{1}(x)+\left(x-x_{0}\right)\left(x-x_{1}\right) \&\left[x_{0}, x_{1}, x_{2}\right] \\
& =\{-0,25587 x+1,40192\}+(x-1,570796) \\
& \times(x-2,0943951) \times(-0,4471007537) \\
& =\{-0,25587 x+1,40192\}+(-0,4471007537) \\
& \times\left\{\begin{array}{cc}
x^{2}-2,0943951 x & +1,570796 \\
-1,570796 x & \times 2,0943951
\end{array}\right\} \\
& =\{-0,25587 x+1,40192\}+\left\{-0,44710 x^{2}+1,63871 x\right. \\
& \text {-1.47090\} } \\
& =-0,44710 x^{2}+1,38284 x-0,06898 \text {. }
\end{aligned}
\]

\[
\begin{aligned}
& P_{3}(x)= P_{2}(x)+\left(x-x_{0}\right)\left(x-x_{1}\right)\left(x-x_{2}\right) f\left[x_{0}, x_{1}, x_{2}, x_{3}\right] \\
&=\left\{\begin{array}{l}
\left.-0,44710 x^{2}+1,38284 x-0,06898\right\} \\
+\{(x-1,57080)(x-2,09439)(x-2,351619)
\end{array}\right. \\
&+\left\{\begin{array}{c}
(0,09126)(x-1,57080) \times \\
\left.\left[x^{2}-2,35619 x+2,09436\right)\right\} \\
-2,09439 x+2,351619] \\
x^{2}-4,45058 x+4,93478 \\
11 \\
+(0,09126 x-0,143351) \times\left(x^{2}-4,45058 x\right. \\
+4,93478)
\end{array}\right. \\
&= 0,09126 x^{3}-0,143351 x^{2}+1,38284 x-0,06898 \\
&-0,40616 x^{2}+0,45035 x-0,707406 \\
&-0,44710 x^{2}+0,637995 x
\end{aligned}
\]

\(P_{3}(x)=0,09126 x^{3}-0,99661 x^{2}+2,471185 x-0,77639\)\\
2. Verification

\begin{itemize}
  \item \(x=\frac{\pi}{2}\)
\end{itemize}

\[
\begin{aligned}
P_{3}(1,570796) & =0,09126 \times 1,570796^{3}=+0,353704 \\
& -0,99661 \times 1,570796^{2}=-2,459035587 \\
& +2,471185 \times 1,570796=+3,881727513 \\
& -0,77639
\end{aligned}
\]

1,000006

\[
\left|P_{3}(x)-\sin x\right|=610^{-6}
\]

\begin{itemize}
  \item \(x=\frac{5 \pi}{6}\)
\end{itemize}

\[
\begin{aligned}
\begin{aligned}
P_{3}(2,617993876) & =0,09126 \times 2,617993876^{3}= \\
& -0,99661 \times 2,617993876^{2}= \\
& +2,471185 \times 2,617993876= \\
& +6,469547196 \\
& =0,77639 \\
& \\
\left|P_{3}(x)-\sin x\right|= & 1,910^{-5}
\end{aligned}
\end{aligned}
\]

0,8660368426\\
\(\left|P_{3}(x)-\sin x\right|=1.1410^{-5}\)

\begin{itemize}
  \item \(x=\frac{3 \pi}{4}\)
\end{itemize}

\[
\begin{aligned}
P_{3}(2,35619449) & =0,09126 \times 2,35619449^{3}=+1,19375341 \\
& -0,99661 \times 2,35619449^{2}=-5,532832373 \\
& +2,471185 \times 2,35619449=+5,822592481 \\
& =0,77639
\end{aligned}
\]

0.707121448

\[
\left|P_{3}(x)-\sin x\right|=1,46710^{-5}
\]

Paur \(x \in\left[\frac{\pi}{2}, \frac{5 \pi}{6}\right]\), l'approximation desin \(x\) ( \(x\) en radians) par le polynôme \(P_{3}(x)\) tel que:

\[
P_{3}(x)=0,09126 x^{3}-0,99661 x^{2}+2,471185 x-0,77639
\]

est satisfaisante. L'erreur absolve\\
\(\left|P_{3}(x)-\sin x\right|\) sur cet intervalle est de l'ordre de \(10^{-5}\)\\
3. Regardons les valeurs de \(P_{3}(x)\) sur \(\left.\left.x \in\right] \frac{5 \pi}{6}, \frac{7 \pi}{6}\right]\), xtoujours exprimé en radians.\\
Calculons parexemple

\[
\begin{aligned}
P_{3}(\pi) L=0,09126 \times 31,00627668= & +2,82963281 \\
\pi=3,141592654-0,99661 \times 9,869604401 & =-9,836146442 \\
+2,171185 \times 3,141592654 & =+7,763456642 \\
& -0,77639 \\
& -0,0194469902
\end{aligned}
\]

\[
\begin{aligned}
\left|P_{3}(x)-\sin x\right| & =0,019446 \\
& \sim 210^{-2}
\end{aligned}
\]

Nous voyons donc pu'en dehors de fintervalle d'interpolation, \(\left[\frac{\pi}{2}, \frac{5 \pi}{6}\right]\), le polynôme \(P_{3}(x)\) n'est pas une bonne approximation de la fonction simus ae -

\section*{Exercice 3:}
\begin{enumerate}
  \item Décomposition LU par la méthode de Doolittle (voir l'exemple du cours)
\end{enumerate}

\[
\begin{aligned}
A & =\left(\begin{array}{cccc}
-2 & 1 & -1 & 1 \\
2 & 0 & 4 & -3 \\
-4 & -1 & -12 & 9 \\
-2 & 1 & 1 & -4
\end{array}\right) \\
& \rightarrow\left(\begin{array}{cccc}
-2 & 1 & -1 & 1 \\
2 & 0 & 4 & -3 \\
-4 & -1 & -12 & 9 \\
-2 & 1 & 1 & -4
\end{array}\right) \rightarrow\left(\begin{array}{cccc}
-2 & 1 & -1 & 1 \\
-1 & 1 & 3 & -2 \\
2 & -3 & -10 & 7 \\
1 & 0 & 2 & -5
\end{array}\right) \rightarrow\left(\begin{array}{cccc}
-2 & 1 & -1 & 1 \\
-1 & 1 & 3 & -2 \\
2 & -3 & -1 & 1 \\
1 & 0 & 2 & -5
\end{array}\right) \\
& \rightarrow\left(\begin{array}{cccc}
-2 & 1 & -1 & 1 \\
-1 & 1 & 3 & -2 \\
2 & -3 & -1 & 1 \\
1 & 0 & -2 & -3
\end{array}\right)
\end{aligned}
\]

On en déduit alors\\
\(L=\left(\begin{array}{cccc}1 & 0 & 0 & 0 \\ -1 & 1 & 0 & 0 \\ 2 & -3 & 1 & 0 \\ 1 & 0 & -2 & 1\end{array}\right)\)

\[
U=\left(\begin{array}{cccc}
-2 & 1 & -1 & 1 \\
0 & 1 & 3 & -2 \\
0 & 0 & -1 & 1 \\
0 & 0 & 0 & -3
\end{array}\right)
\]

\begin{enumerate}
  \setcounter{enumi}{1}
  \item La solution du système \(A x=b\) s'obtient en résolvant successivement les deux systèmes triangulaires\\
\(\left\{\begin{array}{l}L y=b \\ U x=y\end{array}\right.\)\\
\(L y=b \Rightarrow\left(\begin{array}{cccc}1 & 0 & 0 & 0 \\ -1 & 1 & 0 & 0 \\ 2 & -3 & 1 & 0 \\ 1 & 0 & -2 & 1\end{array}\right)\left(\begin{array}{l}y_{1} \\ y_{2} \\ y_{3} \\ y_{4}\end{array}\right)=\left(\begin{array}{c}1.5 \\ 4 \\ -14 \\ -6.5\end{array}\right)\)\\
On a alors \(y=\left(\begin{array}{c}1.5 \\ 5.5 \\ -0.5 \\ -9\end{array}\right)\)\\
\(U x=y \Rightarrow\left(\begin{array}{cccc}-2 & 1 & -1 & 1 \\ 0 & 1 & 3 & -2 \\ 0 & 0 & -1 & 1 \\ 0 & 0 & 0 & -3\end{array}\right)\left(\begin{array}{c}x_{1} \\ x_{2} \\ x_{3} \\ x_{4}\end{array}\right)=\left(\begin{array}{c}1.5 \\ 5.5 \\ -0.5 \\ -9\end{array}\right)\)\\
On obtient \(x=\left(\begin{array}{c}-0.5 \\ 1 \\ 3.5 \\ 3\end{array}\right)\)
\end{enumerate}

\section*{Exercice 4 :}
\begin{enumerate}
  \item Le système s'écrit sous forme matricielle suivante
\end{enumerate}

\[
\left[\begin{array}{ccc}
4 & 2 & 1 \\
-1 & 2 & 0 \\
2 & 1 & 4
\end{array}\right]\left[\begin{array}{l}
x_{1} \\
x_{2} \\
x_{3}
\end{array}\right]=\left[\begin{array}{c}
11 \\
3 \\
16
\end{array}\right]
\]

\section*{Méthode de Jacobi}
La suite récursive de Jacobi permet de calculer la solution à l'itération \(n+1\) en fonction de la solution à l'itération n d'une façon indépendante.

La matrice \(A\) s'écrit de la façon suivante \(A=M-N\), avec

\[
M=\left[\begin{array}{lll}
4 & 0 & 0 \\
0 & 2 & 0 \\
0 & 0 & 4
\end{array}\right] \text { et } \quad \mathrm{N}=\left[\begin{array}{ccc}
0 & -2 & -1 \\
1 & 0 & 0 \\
-2 & -1 & 0
\end{array}\right]
\]

On a alors

\[
A x=b \Leftrightarrow(M-N) x=b
\]

On obtient la suite récursive de Jacobi suivante

\[
M x^{n+1}=N x^{n}+b
\]

\[
\begin{aligned}
& {\left[\begin{array}{lll}
4 & 0 & 0 \\
0 & 2 & 0 \\
0 & 0 & 4
\end{array}\right]\left[\begin{array}{l}
x_{1}^{n+1} \\
x_{2}^{n+1} \\
x_{3}^{n+1}
\end{array}\right]=\left[\begin{array}{ccc}
0 & -2 & -1 \\
1 & 0 & 0 \\
-2 & -1 & 0
\end{array}\right]\left[\begin{array}{l}
x_{1}^{n} \\
x_{2}^{n} \\
x_{3}^{n}
\end{array}\right]+\left[\begin{array}{c}
11 \\
3 \\
16
\end{array}\right]} \\
& {\left[\begin{array}{l}
x_{1}^{n+1} \\
x_{2}^{n+1} \\
x_{3}^{n+1}
\end{array}\right]=\left[\begin{array}{ccc}
\frac{1}{4} & 0 & 0 \\
0 & \frac{1}{2} & 0 \\
0 & 0 & \frac{1}{4}
\end{array}\right]\left[\begin{array}{ccc}
0 & -2 & -1 \\
1 & 0 & 0 \\
-2 & -1 & 0
\end{array}\right]\left[\begin{array}{l}
x_{1}^{n} \\
x_{2}^{n} \\
x_{3}^{n}
\end{array}\right]+\left[\begin{array}{ccc}
\frac{1}{4} & 0 & 0 \\
0 & \frac{1}{2} & 0 \\
0 & 0 & \frac{1}{4}
\end{array}\right]\left[\begin{array}{c}
11 \\
3 \\
16
\end{array}\right]} \\
& {\left[\begin{array}{l}
x_{1}^{n+1} \\
x_{2}^{n+1} \\
x_{3}^{n+1}
\end{array}\right]=\left[\begin{array}{ccc}
0 & -\frac{1}{2} & -\frac{1}{4} \\
\frac{1}{2} & 0 & 0 \\
-\frac{1}{2} & -\frac{1}{4} & 0
\end{array}\right]\left[\begin{array}{l}
x_{1}^{n} \\
x_{2}^{n} \\
x_{3}^{n}
\end{array}\right]+\left[\begin{array}{c}
11 / 4 \\
3 / 2 \\
4
\end{array}\right]}
\end{aligned}
\]

Prenons la valeur de départ ( \(n=0\) ) suivante :

\[
\left[\begin{array}{l}
x_{1}^{0} \\
x_{2}^{0} \\
x_{3}^{0}
\end{array}\right]=\left[\begin{array}{l}
0 \\
0 \\
0
\end{array}\right]
\]

Après 4 itérations, on obtient :

\[
\begin{aligned}
& n=1 \\
& {\left[\begin{array}{l}
x_{1}^{1} \\
x_{2}^{1} \\
x_{3}^{1}
\end{array}\right]=\left[\begin{array}{c}
2,75 \\
1,5 \\
4
\end{array}\right]} \\
& n=2 \\
& {\left[\begin{array}{l}
x_{1}^{2} \\
x_{2}^{2} \\
x_{3}^{2}
\end{array}\right]=\left[\begin{array}{c}
1 \\
2,875 \\
2.25
\end{array}\right]} \\
& n=3 \\
& {\left[\begin{array}{l}
x_{1}^{3} \\
x_{2}^{3} \\
x_{3}^{3}
\end{array}\right]=\left[\begin{array}{c}
0,75 \\
2 \\
2,78125
\end{array}\right]} \\
& n=4
\end{aligned}
\]

\[
\left[\begin{array}{l}
x_{1}^{4} \\
x_{2}^{4} \\
x_{3}^{4}
\end{array}\right]=\left[\begin{array}{c}
1.0546875 \\
1.875 \\
3.125
\end{array}\right]
\]

Vérification que la matrice A est à diagonale strictement dominante

\[
\begin{aligned}
\sum_{j=1}^{N}\left|a_{i j}\right|<\left|a_{i i}\right| \quad \text { avec } \mathrm{i} & =1, \ldots, \mathrm{~N} \\
\mathrm{j} & =1, \ldots, \mathrm{~N} \text { et } \mathrm{j} \neq \mathrm{i}
\end{aligned}
\]

On a

\[
\begin{array}{ll}
i=1 & \sum_{j \neq i}^{N}\left|a_{i j}\right|=2+1=3<a_{11}=4 \\
i=2 & \sum_{j \neq i}^{N}\left|a_{i j}\right|=+1<a_{22}=2 \\
i=3 & \sum_{j \neq i}^{N}\left|a_{i j}\right|=2+1=3<a_{33}=4
\end{array}
\]

La matrice A est bien à diagonale strictement dominante

\section*{Méthode de Gauss-Seidel}
La suite récursive de Gauss-Seidel permet de calculer la composante J de la solution à l'itération \(n+1\) à partir des composantes \(j-1, j-2, \ldots, 1\) disponibles à l'itération \(n+1\) et de la solution à l'itération \(n\).\\
On a les matrices \(M\) et \(N\) suivantes :

\[
M=\left[\begin{array}{ccc}
4 & 0 & 0 \\
-1 & 2 & 0 \\
2 & 1 & 4
\end{array}\right] \text { et } \quad \mathrm{N}=\left[\begin{array}{ccc}
0 & -2 & -1 \\
0 & 0 & 0 \\
0 & 0 & 0
\end{array}\right]
\]

On a alors

\[
A x=b \Leftrightarrow(M-N) x=b
\]

On obtient la suite récursive de Gauss-Seidel suivante \(M x^{n+1}=N x^{n}+b\)

\[
\begin{aligned}
& {\left[\begin{array}{ccc}
4 & 0 & 0 \\
-1 & 2 & 0 \\
2 & -1 & 4
\end{array}\right]\left[\begin{array}{l}
x_{1}^{n+1} \\
x_{2}^{n+1} \\
x_{3}^{n+1}
\end{array}\right]=\left[\begin{array}{ccc}
0 & -2 & -1 \\
0 & 0 & 0 \\
0 & 0 & 0
\end{array}\right]\left[\begin{array}{l}
x_{1}^{n} \\
x_{2}^{n} \\
x_{3}^{n}
\end{array}\right]+\left[\begin{array}{c}
11 \\
3 \\
16
\end{array}\right]} \\
& \left\{\begin{array}{l}
x_{1}^{n+1}=11 / 4-1 / 2 x_{2}^{n}-1 / 4 x_{3}^{n} \\
x_{2}^{n+1}=3 / 2+1 / 2 x_{1}^{n+1} \\
x_{3}^{n+1}=4-1 / 2 x_{1}^{n+1}-1 / 4 x_{2}^{n+1}
\end{array}\right.
\end{aligned}
\]

Prenons la valeur de départ ( \(n=0\) ) suivante \(\left[\begin{array}{l}x_{1}^{0} \\ x_{2}^{0} \\ x_{3}^{0}\end{array}\right]=\left[\begin{array}{l}0 \\ 0 \\ 0\end{array}\right]\)

Après 4 itérations, on obtient

\[
\begin{aligned}
& n=1 \\
& {\left[\begin{array}{l}
x_{1}^{1} \\
x_{2}^{1} \\
x_{3}^{1}
\end{array}\right]=\left[\begin{array}{c}
2,75 \\
2,875 \\
1,90625
\end{array}\right]} \\
& n=2 \\
& {\left[\begin{array}{l}
x_{1}^{2} \\
x_{2}^{2} \\
x_{3}^{2}
\end{array}\right]=\left[\begin{array}{c}
0,835975 \\
1,9179875 \\
3,102515625
\end{array}\right]} \\
& n=3 \\
& {\left[\begin{array}{l}
x_{1}^{3} \\
x_{2}^{3} \\
x_{3}^{3}
\end{array}\right]=\left[\begin{array}{c}
1,015377344 \\
2,0076887672 \\
2,99038916
\end{array}\right]}
\end{aligned}
\]


\end{document}